\chapter{Conclusion and Future Work}
\label{ch:conclusion}
This chapter draws together the work of FYP-I. It summarises what was done and what was found, states which objectives were met and which were not, records the main challenges, and sets out the plan of work for FYP-II. Because the real cohort did not arrive during this phase, every result summarised here comes from a simulation in which the true relationship between genotype and face is known, and the chapter keeps that distinction in view throughout: what the project can claim is that its pipeline and its evaluation behave correctly, not that genotype predicts real faces to any particular accuracy. The plan at the end is therefore built around the one dependency the team does not control, which is the release of the data.

\section{Summary of the Work}
FaceCon set out to build a reproducible pipeline that predicts three-dimensional facial shape from genotype and reports honestly how much the genotype adds beyond age, sex, and ancestry. The earlier chapters established why such a pipeline is needed, what it must do, and how it is designed. During this phase the pipeline was implemented from the genotype matrix through to a decoded face and a set of evaluation measures, together with an integrated appearance module for eye, hair, and skin colour.

Access to the real cohort was not granted in time, so the pipeline could not be run on the data it was designed for. The team used this period to build a simulated cohort that follows the public description of the target study, to add published association signals to it, and to test the pipeline where the correct answer is known. The simulation was then made substantially harder by replacing its generated genotypes with real human genotypes from the 1000 Genomes Project, which supply real correlation between SNPs, real allele frequencies, and continuous population structure, while the faces remained simulated. This turned a delay into a validation exercise that would have been needed in any case.

\section{Findings}
The findings of this phase are about the pipeline and not about human faces. The pipeline recovers a planted genetic contribution at its true size and reports no contribution under two independent negative controls. It recovers all of the published associations that were placed in the simulated measurements and reports no others. These properties hold on real genotypes as well as on generated ones, and a set of 44 automated tests guards them.

The experiments also produced three findings about method. The first is that the way SNPs are presented to a model can matter more than the model itself. Compressing SNPs to PCs removed almost all of the planted signal, and led us to a provisional conclusion that was wrong until a positive control exposed it. The second is that, on an additive planted signal and a cohort of a few thousand subjects, a regularised linear model outperformed every neural model, by $0.077$ in $R^2$ on real genotypes where both were given identical inputs on one split, and more data narrowed that gap without closing it. The third comes from the real genotypes and could not have come from the simulation. When the causal SNPs are hidden, so that models can use only the SNPs correlated with them, about a fifth of the attainable signal survives. This is a mechanism for the small effects that published association studies report, since a panel of common SNPs reads neighbours rather than causes.

A fourth finding came from deliberately planting the non-additive structure that the simulation had so far omitted. When a third of the genetic share is carried by heterozygote deviations and a third by interactions between pairs of SNPs, a model given the true causal SNPs as additive dosages recovers $0.058$ in $R^2$, while the same model given the heterozygote codes and the true interacting pairs recovers $0.148$. The recoverable total is nearly unchanged across every architecture tested, so the shortfall is not missing signal but signal in a form the additive model cannot express. This sharpens the first finding: representation dominates architecture for non-additive signal as well as additive. Running the neural models on those same cohorts settled the question the additive comparison had to leave open. With the linear model's home advantage removed, the gap did not close but widened, to $-0.065$ with a $t$ interval of $[-0.070,-0.061]$ across three cohorts, and the transformer scored the same on permuted genotypes as on real ones, meaning it drew nothing from them at all. The claim this licenses is narrow and is stated narrowly: at roughly 2,200 training subjects, these architectures cannot exploit structure a regularised linear model recovers from the same features, and scale rather than architecture is the first thing a successor should change. It also produced a negative result worth stating, in that the practical two-degree-of-freedom screen helps on dominance and actively hurts on pure epistasis, because a per-SNP score cannot see a pair whose individual effects average out. All four findings will shape how the real data is analysed, and none should be read as a claim about real facial genetics.

\section{Objectives Met and Not Met}
Table~\ref{tab:objectives} sets the objectives of Section~\ref{sec:goals} against what was achieved.

\begin{table}[!ht]
\centering
\caption{Status of the project objectives at the end of FYP-I.}
\label{tab:objectives}
\small\setstretch{1.15}
\begin{tabular}{|P{4.9cm}|P{2.2cm}|P{7.1cm}|}
\hline
\textbf{Objective} & \textbf{Status} & \textbf{Remark} \\ \hline
Reproducible genotype to face pipeline on paired 3DFN subjects & Partly met & Built and validated on simulated data. Not yet run on 3DFN subjects \\ \hline
Documented genotype QC and SNP panel & Partly met & The filtering, relatedness removal, and allele alignment steps are implemented, documented, and exercised on the 1000 Genomes files. They have not been run on the 3DFN genotypes, which have not arrived. The panel documentation schema is defined in Appendix~A \\ \hline
Documented scan to FLAME fitting & Partly met & The fitting procedure recovers the shape of synthetic scans to within $0.5$~mm and tells a face from a random mesh, as Section~\ref{sec:res-flame} reports. It has not been run on real scans, and it fits the stand-in face model rather than FLAME \\ \hline
Progressively more complex models with baselines & Met & Three baselines and four neural models share one interface and one evaluation \\ \hline
Rigorous evaluation with ablations & Partly met & Parameter, mesh, regional, and landmark measures, controls, a power analysis, an ablation that hides the causal SNPs, and non-additive cohorts are complete. The comparison of published and random SNP panels for RQ4 was attempted and proved unable to detect the planted effects, so it is reported as a limitation \\ \hline
Separation of genetic from demographic signal & Met in simulation & Demonstrated with positive and negative controls \\ \hline
Appearance through an established method & Partly met & Integrated, and reproduces the official calculator. A check against self-reported eye colour is implemented, but no open cohort with that information remains available, as Section~\ref{sec:res-appearance-untested} explains \\ \hline
Limitations, ethics, and a demonstration & Partly met & Limitations are stated throughout. The demonstrator is built from the corrected pipeline and shows the results of Chapter~\ref{ch:results}, but only for simulated faces until the real data arrives. The presentation slides still predate the corrections \\ \hline
\end{tabular}
\end{table}

The scope of the project was therefore not entirely covered. Everything that could be done without the real data was done, and everything that depends on the real data remains. The objectives that are only partly met share a single cause: each was completed and validated as far as the simulated and public data allow, and each now waits on the release of the controlled-access cohort to be finished on the data it was designed for.

\section{Challenges}
The main challenge was access to data. Both parts of the cohort are released only after an application with ethical approval is reviewed by a data access committee, and that process did not finish within FYP-I. The second challenge was the correctness of our own analysis. A review during this phase found six defects, two of which had affected reported numbers. They are listed in Chapter~\ref{ch:impl}, and the lesson taken from them is that every result must be accompanied by a control that could have contradicted it. The third challenge was computing resources, since the graphics card was available for limited periods, which restricted the neural experiments to three training runs per model.

\section{Recommendations for Future Work}
Beyond the plan for FYP-II, three directions would strengthen work in this area. The first is evaluation on cohorts of other ancestries, since both the target cohort and most published associations are European. The second is the study of non-linear and interaction effects. Section~\ref{sec:res-nonadd} takes a first step by planting them and measuring what an additively restricted model loses, but it leaves the harder problem open: a screen that can find interacting pairs rather than being told them, which must be specified before it is run so that its result means something, and the question of scale, since the neural comparison on those cohorts has now been run and lost at this sample size. The third is a common public benchmark with fixed splits and mandatory demographic baselines, so that results from different groups can be compared.

\section{Plan of Work for FYP-II}
The plan has two tracks, because the arrival of the real data is outside the team's control. Track~A is the work that does not depend on the real data and will be carried out in any case. Track~B begins when access is granted. If access is not granted by the middle of the semester, Track~A continues to completion, and the project will be submitted as a validated pipeline with a documented protocol for real data. Table~\ref{tab:fyp2} gives the schedule in weeks of the semester.

\begin{table}[!ht]
\centering
\caption{Plan of work for FYP-II.}
\label{tab:fyp2}
\small\setstretch{1.15}
\begin{tabular}{|P{2.3cm}|P{1.3cm}|P{10.6cm}|}
\hline
\textbf{Weeks} & \textbf{Track} & \textbf{Work} \\ \hline
1 to 2 & A & Correct the presentation slides so that they show the results of Chapter~\ref{ch:results}, as the demonstrator already does \\ \hline
2 to 5 & A & Specify in advance, and then run, a screen that can find interacting pairs of SNPs rather than being told them, the open problem left by Section~\ref{sec:res-nonadd} \\ \hline
3 to 6 & A & Test whether the failed positive control on the real-genotype mixed cohort is a matter of power, by repeating it with fewer causal SNPs or a larger training split, as the protocol records \\ \hline
5 to 8 & A & Replace the stand-in face model with FLAME once its licence is accepted, and repeat the scan-fitting validation on it \\ \hline
from access & B & Apply genotype QC, document the SNP panel, register the scans to FLAME, and pair the subjects \\ \hline
after B begins & B & Run the fixed protocol on the real data without change: the baseline ladder, the controls, the comparison of models, and the regional and subgroup analyses \\ \hline
12 to 14 & A and B & Update the user manual for the real-data workflow, regenerate the demonstrator from the real-data results, and prepare the final report and the code release \\ \hline
15 to 16 & A and B & Final evaluation, defence preparation, and submission \\ \hline
\end{tabular}
\end{table}

Several items that earlier versions of this plan scheduled for FYP-II were completed during FYP-I and are reported in Chapter~\ref{ch:results}: the experiment on real genotypes from the 1000 Genomes Project, the cohorts with non-additive effects and the neural comparison on them, the validation of scan fitting on synthetic scans, the regional and subgroup analyses, and repeated cohorts with confidence intervals for the neural models. Two were attempted and could not be completed. The check of the colour component against self-reported eye colour has no open cohort left to run on, and the comparison of published and random SNP panels proved unable to detect the effects it was designed to find. Both are recorded as limitations rather than carried forward, because neither can be repaired by more work within the project.

The analysis protocol for the real data has been written and frozen, together with the thresholds by which each result will be judged, so that nothing can be adjusted after the real results are seen. Three of its decisions matter most. The primary measure is the gain in test $R^2$ from adding SNPs to covariates and ancestry components. A result will be reported as evidence of genetic signal only if its confidence interval lies above zero and the shuffled-genotype control does not. A finding that SNPs add little will be reported as plainly as a positive one.

\section{Conclusion}
FYP-I delivered a working and tested pipeline for predicting facial shape from genotype, an integrated appearance module, and a body of evidence that the pipeline measures genetic signal correctly and does not invent it. It did not deliver results on real faces, because the real data did not arrive. The project is honest about that boundary: what has been shown is that the method is sound and ready, and what remains is to apply it. FYP-II will complete the parts that do not depend on the data, test whether the ranking of models changes when the planted effects are not additive, and run a protocol fixed in advance on the real cohort as soon as access is granted.
